% 4-6-shap.tex
%  Budget: 8pp -- the largest section of chapter 4.  HANDOFF 6.7, 7/RQ3, 19.1
%  (shap_importance.csv), 19.2 (figures 10-15).
%  Source map: lundberg_unified_2017 for the method; trebbien_explainable_2023
%   as the EPF precedent. All numbers from shap_importance.csv.
%  MANDATE:
%   - state the fit boundary EXPLICITLY: the explained model is refit on the
%     trailing 1092 days ending 2017-12-31, NOT the walk-forward artifact
%   - SHAP is model-specific: attribution of one model, not causation
%   - do not describe figure content beyond what the figure inventory states
%  Rules: decimals in \lr{}; integers Persian; every table cell in \lr{}.

\section{تحلیل تفسیرپذیری \lr{SHAP}}
\label{sec:4-6}

\subsection*{پرسش، و آنچه این بخش می‌تواند و نمی‌تواند پاسخ دهد}

بخش‌های پیشین نشان دادند کدام مدل چقدر دقیق است. پرسشِ اینجا از جنس دیگری است:
مدل بر چه چیزی تکیه می‌کند، و آیا این تکیه‌گاه با تغییر شرایط بازار جابه‌جا
می‌شود؟

روش، مقادیر \lr{SHAP} است \cite{lundberg_unified_2017}: چارچوبی که سهم هر ویژگی
در فاصله‌ی میان پیش‌بینیِ یک نمونه و میانگین پیش‌بینی‌ها را به‌صورت جمعی تجزیه
می‌کند. مزیت آن نسبت به شاخص‌های اهمیتِ سنتیِ درختی این است که سهم‌ها جمع‌شونده و
نمونه‌به‌نمونه‌اند، پس می‌توان آن‌ها را روی هر زیرمجموعه‌ای از روزها میانگین گرفت و
مقایسه کرد \lr{—} و کل این بخش دقیقاً بر همین خاصیت بنا شده است. کاربرد آن در
پیش‌بینی قیمت برق سابقه دارد \cite{trebbien_explainable_2023}.

پیش از هر عددی، سه قید دامنه باید ثبت شود.

\textbf{یکم: این تحلیل مختص \lr{LightGBM} است.} پیاده‌سازی از
\lr{TreeExplainer} استفاده می‌کند که برای مدل‌های درختی، مقادیر \lr{SHAP} را
دقیق و نه تقریبی محاسبه می‌کند. چهار مدل دیگر این پژوهش تحلیل نمی‌شوند، و هیچ
نتیجه‌ای از این بخش به آن‌ها تعمیم داده نمی‌شود. \lr{LEAR-LASSO} سازوکار
تفسیرپذیریِ خودش را دارد \lr{—} ضرایبی که بخشی‌شان صفر شده‌اند \lr{—} و
\lr{LSTM} هیچ‌کدام را ندارد. اینکه سه مدل با دقت آماریِ تفکیک‌ناپذیر
(\cref{sec:4-5-1}) از اطلاعات یکسانی استفاده می‌کنند یا نه، پرسشی است که این
پژوهش به آن پاسخ نمی‌دهد.

\textbf{دوم، و مهم‌تر: مدلی که توضیح داده می‌شود، همان شیئی نیست که نتایج از آن
گزارش شده است.} این نکته باید صریح نوشته شود، چون در غیر این صورت خواننده‌ای که
خودش کشفش کند آن را نقص می‌بیند. مدل‌های اعتبارسنجی پیش‌رونده که تمام اعداد
\cref{sec:4-2} تا \cref{sec:4-5-2} از آن‌ها می‌آید، هر یک تنها تا پیش از مبدأ خود
آموزش دیده‌اند. توضیح‌دادنِ پیش‌بینیِ آن مدل‌ها روی همان روزهایی که در
آموزششان بوده‌اند، توضیحی درون‌نمونه‌ای می‌بود. بنابراین برای این بخش یک برازش
جداگانه انجام شده است: \lr{LightGBM} بر پنجره‌ی متحرکِ ۱۰۹۲ روزه‌ای که در
۳۱ دسامبر ۲۰۱۷ پایان می‌یابد، با همان ابرپارامترهای تنظیم‌شده و همان بذر ۴۲.

این برازش، همزادِ وفادارِ مدل گزارش‌شده است و نه خودِ آن. سه چیز میان آن دو یکسان
است: پیکربندی ابرپارامترها، ماتریس ویژگی، و طول پنجره‌ی آموزش. آنچه متفاوت است،
مبدأ پایانیِ ثابت به‌جای مبدأ متحرک است. ماژول تحلیل عامدانه از برازش بر پنجره‌ی
کوتاه‌تر امتناع می‌کند، دقیقاً برای آنکه این همزاد نتواند بی‌سروصدا از مدل اصلی
فاصله بگیرد. \cref{sec:3-5} این را به‌عنوان سومین شیءِ مدلیِ پروژه تفکیک کرده و
گفته است که هیچ عددی از آن وارد جدول‌های دقت نمی‌شود؛ اینجا نیز چنین است. دلیلش
صریح است: پنجره‌ی آموزشِ این همزاد روزهای دوره‌ی آزمون را نیز در بر
می‌گیرد، پس برای سنجش دقت معتبر نیست و تنها برای تفسیر به کار می‌رود. مدل‌های
گزارش‌شده در \cref{sec:4-2} تا \cref{sec:4-5-2}، برعکس، هیچ‌یک فراتر از مبدأ خود
آموزش ندیده‌اند و قاعده‌ی نشت را نقض نمی‌کنند.

\textbf{سوم: \lr{SHAP} نسبت‌دادن است، نه علیت.} آنچه اندازه گرفته می‌شود، سهم یک
ویژگی در خروجیِ \emph{یک مدل} است. اگر مدل بر ویژگی‌ای تکیه کند که تنها همبسته‌ی
علت واقعی است، \lr{SHAP} همان ویژگیِ همبسته را برجسته نشان می‌دهد. هیچ جمله‌ای در
سرتاسر این بخش نباید چنین خوانده شود که «فلان متغیر قیمت برق آلمان را تعیین می‌کند».

\subsection*{تصویر کلی}

\cref{tab:shap-importance} میانگینِ قدرمطلقِ مقادیر \lr{SHAP} را به تفکیک
خانواده‌ی ویژگی و در چهار حالت گزارش می‌کند. واحد، یورو بر مگاوات‌ساعت است، و
تفسیر مستقیم دارد: میانگینِ اندازه‌ی جابه‌جایی که آن خانواده در پیش‌بینی ایجاد
می‌کند.

\begin{table}[htbp]
\centering
\caption{میانگین قدرمطلق مقادیر \lr{SHAP} به تفکیک خانواده‌ی ویژگی، بر حسب یورو
بر مگاوات‌ساعت، برای مدل \lr{LightGBM}. \lr{exog\_1} پیش‌بینی بار و
\lr{exog\_2} پیش‌بینی تولید تجدیدپذیر است؛ پسوند \lr{D0} یعنی مقدارِ خودِ روز
هدف و \lr{D-n} یعنی مقدار \lr{n} روز پیش. مقادیر پررنگ، بیشترین مقدار هر ستون
است.}
\label{tab:shap-importance}
\small
\setlength{\tabcolsep}{5pt}
\begin{tabular}{lrrrr}
\toprule
& \multicolumn{2}{c}{ساعتی} & \multicolumn{2}{c}{روزانه} \\
\cmidrule(lr){2-3}\cmidrule(lr){4-5}
خانواده‌ی ویژگی & آرام & پرتنش & آرام & پرتنش \\
\midrule
\lr{price\_D-1} & \lr{3.4833} & \lr{6.1471} & \lr{3.1109} & \lr{\textbf{5.1335}} \\
\lr{price\_D-2} & \lr{0.8610} & \lr{1.2550} & \lr{0.3778} & \lr{0.5177} \\
\lr{price\_D-3} & \lr{0.8279} & \lr{1.0302} & \lr{0.3328} & \lr{0.3916} \\
\lr{price\_D-7} & \lr{1.7454} & \lr{2.1473} & \lr{1.0802} & \lr{1.4784} \\
\midrule
\lr{exog\_1\_D-1} & \lr{0.6591} & \lr{0.6735} & \lr{0.2479} & \lr{0.2611} \\
\lr{exog\_1\_D-7} & \lr{0.7131} & \lr{0.7487} & \lr{0.3699} & \lr{0.3991} \\
\lr{exog\_1\_D0} & \lr{3.5197} & \lr{3.3680} & \lr{\textbf{3.4393}} & \lr{3.0130} \\
\midrule
\lr{exog\_2\_D-1} & \lr{0.9947} & \lr{0.9221} & \lr{0.4077} & \lr{0.3986} \\
\lr{exog\_2\_D-7} & \lr{0.6740} & \lr{0.7761} & \lr{0.2230} & \lr{0.2754} \\
\lr{exog\_2\_D0} & \lr{\textbf{5.8422}} & \lr{\textbf{5.9761}} & \lr{4.4929} & \lr{5.0366} \\
\midrule
\lr{dow} & \lr{0.1464} & \lr{0.1358} & \lr{0.0395} & \lr{0.0379} \\
\bottomrule
\end{tabular}
\end{table}

\begin{figure}[htbp]
  \centering
  \includegraphics[width=0.85\textwidth]{10_shap_global_importance_hourly}
  \caption{اهمیت سراسری ویژگی‌ها بر پایه‌ی میانگین قدرمطلق مقادیر \lr{SHAP}،
  افق ساعتی.}
  \label{fig:shap-global}
\end{figure}

\begin{figure}[htbp]
  \centering
  \includegraphics[width=0.85\textwidth]{11_shap_beeswarm_hourly}
  \caption{نمودار زنبوری (\lr{beeswarm}) مقادیر \lr{SHAP} در افق ساعتی: پراکندگی سهم هر ویژگی در
  تک‌تک نمونه‌ها، نه تنها میانگین آن.}
  \label{fig:shap-beeswarm}
\end{figure}

سه واقعیت از \cref{tab:shap-importance} و \cref{fig:shap-global} بیرون می‌آید.

\textbf{پیش‌بینی تجدیدپذیرِ روز هدف، بزرگ‌ترین سهم را دارد.}
\lr{exog\_2\_D0} با \lr{5.8422} در حالت آرام و \lr{5.9761} در حالت پرتنش، در هر
دو حالتِ افق ساعتی بالاترین است. سازوکار بازار همین را انتظار دارد: تولید بادی و
خورشیدیِ پیش‌بینی‌شده، منحنی عرضه‌ی باقی‌مانده را جابه‌جا می‌کند و مستقیماً بر
قیمت تسویه اثر می‌گذارد. اینکه مدل چنین وزنی به آن می‌دهد، نشانه‌ی خوبی است که
چیزی واقعی را آموخته و صرفاً پایداریِ سری را دنبال نمی‌کند.

\textbf{دو ویژگیِ روز هدف، بیشتر از هر لگی وزن دارند.} \lr{exog\_2\_D0} و
\lr{exog\_1\_D0} روی هم بخش بزرگی از کل اهمیت را می‌گیرند. همان استدلالی که
\cref{sec:3-4} برای مجاز شمردنِ ویژگی‌های \lr{D0} آورد، اینجا هم برقرار است: پیش‌بینی بار و
تجدیدپذیرِ روز هدف، پیش از بسته‌شدن مزایده در دسترس شرکت‌کنندگان است، پس
استفاده از آن‌ها نشت اطلاعات نیست. اگر این ویژگی‌ها کنار گذاشته می‌شدند، مدل از
دو مؤثرترین منبع اطلاعاتی خود محروم می‌ماند.

\textbf{روز هفته تقریباً بی‌اثر است.} \lr{dow} با \lr{0.1464} در افق ساعتی و
\lr{0.0395} در افق روزانه، حدود یک‌چهلم بزرگ‌ترین خانواده است. خواندنِ آن به معنای
نبودِ الگوی هفتگی نادرست است \lr{—} \cref{sec:3-3-3} نشان داد که چنین الگویی هست
\lr{—} بلکه به این معناست که لگ هفتگیِ قیمت، \lr{price\_D-7} با \lr{1.7454}،
همان اطلاعات را در شکلی غنی‌تر منتقل می‌کند. وقتی لگ‌ها حاضرند، متغیر تقویمی
چیزی برای افزودن ندارد.

\subsection*{نیم‌رخ ساعتی}

\begin{figure}[htbp]
  \centering
  \includegraphics[width=0.85\textwidth]{12_shap_hour_profile}
  \caption{تفکیک سهم \lr{SHAP} بر حسب ساعت شبانه‌روز.}
  \label{fig:shap-hour}
\end{figure}

\cref{fig:shap-hour} همان تجزیه را بر محور ساعت شبانه‌روز باز می‌کند. تفکیکِ ساعتی
یادآور می‌شود که میانگین‌های \cref{tab:shap-importance} روی ۲۴ ساعتِ ناهمگن گرفته
شده‌اند: ساعت‌های میانه‌ی روز، که تولید خورشیدی در آن‌ها متمرکز است، رفتاری
متفاوت از ساعت‌های شب دارند، و \cref{sec:3-3-3} همین ناهمگنی را در سطح خودِ قیمت
نشان داده بود (کمینه‌ی میانگین در ساعت ۳ و بیشینه در ساعت ۱۹).

\subsection*{آرام در برابر پرتنش: یافته‌ی اصلی این بخش}

\begin{figure}[htbp]
  \centering
  \includegraphics[width=0.85\textwidth]{13_shap_calm_vs_stressed}
  \caption{مقایسه‌ی اهمیت ویژگی‌ها میان روزهای آرام و روزهای پرتنش.}
  \label{fig:shap-regime}
\end{figure}

مقایسه‌ی دو ستون نخستِ \cref{tab:shap-importance} روشن‌ترین یافته‌ی این بخش را
می‌دهد، و \cref{fig:shap-regime} همان را به‌صورت تصویری نشان می‌دهد.

\textbf{اهمیت \lr{price\_D-1} زیر تنش \lr{76.5} درصد افزایش می‌یابد}، از \lr{3.4833}
به \lr{6.1471}. هیچ خانواده‌ی دیگری چنین جابه‌جایی‌ای ندارد. وقتی بازار پرتنش
است، مدل به‌مراتب سنگین‌تر بر قیمت دیروز تکیه می‌کند.

در جهت مخالف، \lr{exog\_1\_D0} \lr{—} پیش‌بینی بار روز هدف \lr{—} از \lr{3.5197}
به \lr{3.3680} \emph{کاهش} می‌یابد، یعنی حدود چهار درصد افت. و
\lr{exog\_2\_D0} با \lr{5.8422} در برابر \lr{5.9761} عملاً ثابت می‌ماند.

خوانشِ این سه عدد با هم، سازوکار روشنی می‌سازد. در بازار آرام، قیمت را می‌توان
تا حد خوبی از بنیان‌های همان روز \lr{—} بار و تجدیدپذیر \lr{—} ساخت. وقتی
بازار پرتنش می‌شود، همان بنیان‌ها دیگر برای تعیین سطح قیمت کافی نیستند، چون
قیمت‌های بالا از کمیابی و رفتار پیشنهاددهی می‌آیند که در متغیرهای بنیادی
پیدا نیست. مدل در پاسخ، سطح را از آخرین مشاهده‌ی محقق‌شده می‌گیرد و بنیان‌ها را
برای تعیین شکلِ درون‌روزی نگه می‌دارد.

سایر لگ‌های قیمت نیز در همین جهت حرکت می‌کنند اما با شدت کمتر:
\lr{price\_D-2} از \lr{0.8610} به \lr{1.2550} و \lr{price\_D-7} از \lr{1.7454}
به \lr{2.1473}. پس این افزایشِ اتکا، مختص یک لگ نیست بلکه به کل بلوک لگ‌های
قیمت مربوط است، با بیشترین سهم برای نزدیک‌ترین آن‌ها.

یافته‌ی بالا با نتیجه‌ی \cref{sec:4-5-1} پیوند مستقیم دارد و آن را از سوی دیگری
تأیید می‌کند. آنجا نشان داده شد که تمام سود وزن‌دهیِ رژیم‌محور از همان ۷۷ روز
پرتنش می‌آید. اینجا دیده می‌شود که چرا: در آن روزها نه فقط خطای مدل، بلکه
رفتارش واقعاً تغییر می‌کند. دو تحلیلِ کاملاً مستقل \lr{—} یکی بر پایه‌ی زیانِ
پیش‌بینی و دیگری بر پایه‌ی نسبت‌دادنِ ویژگی \lr{—} به یک نتیجه می‌رسند. چنین
هم‌گرایی‌ای خودْ شاهدی است بر اینکه تفکیک دو رژیم صرفاً یک برشِ دلبخواهِ داده
نبوده است.

\subsection*{ساعتی در برابر روزانه}

\begin{figure}[htbp]
  \centering
  \includegraphics[width=0.85\textwidth]{14_shap_hourly_vs_daily}
  \caption{مقایسه‌ی اهمیت ویژگی‌ها میان افق ساعتی و افق روزانه.}
  \label{fig:shap-horizon}
\end{figure}

دو ستون آخر \cref{tab:shap-importance} همان تحلیل را برای هدف روزانه تکرار
می‌کنند، و الگویی متفاوت اما سازگار پدید می‌آورند.

بیشتر اهمیت‌ها در افق روزانه کوچک‌تر می‌شوند، که انتظار می‌رفت: هدف روزانه
میانگینِ ۲۴ عدد است و دامنه‌ی نوسان آن کمتر، پس اندازه‌ی جابه‌جایی‌ها هم کمتر
است. آنچه معنا دارد، نه اندازه‌ی مطلق بلکه نسبت‌هاست.

دو تغییرِ نسبی برجسته است. نخست، لگ‌های میانی به‌شدت افت می‌کنند:
\lr{price\_D-2} از \lr{0.8610} به \lr{0.3778} و \lr{price\_D-3} از \lr{0.8279}
به \lr{0.3328}، یعنی بیش از نصف. دوم، \lr{exog\_1\_D0} تقریباً دست‌نخورده
می‌ماند: \lr{3.5197} در افق ساعتی در برابر \lr{3.4393} در افق روزانه. در واقع
پیش‌بینی بار در افق روزانه به بزرگ‌ترین ویژگی در حالت آرام بدل می‌شود.

توضیح، در ماهیت خودِ دو هدف است. بخش بزرگی از کاری که لگ‌های قیمت در افق ساعتی
انجام می‌دهند، بازسازیِ شکلِ درون‌روزی است \lr{—} اینکه ساعت ۱۹ گران‌تر از ساعت
۳ است. در هدف روزانه این شکل با میانگین‌گیری از میان می‌رود و آن بخش از کارکرد
لگ‌ها بی‌مصرف می‌شود. آنچه باقی می‌ماند، تعیینِ سطح است، که بار و تجدیدپذیرِ
تجمیع‌شده‌ی همان روز آن را به‌خوبی توضیح می‌دهند.

تصویر بالا با یافته‌ی \cref{sec:4-4} کاملاً هم‌خوان است. آنجا دیده شد که برای سه مدل از
چهار مدل، میانگین‌گیری از ۲۴ پیش‌بینی ساعتی بهتر از مدل‌سازی مستقیمِ هدف روزانه
است. اینجا سازوکارش دیده می‌شود: مسیر مستقیم، مدل را از اطلاعاتِ ساختار
درون‌روزی محروم می‌کند و همان اطلاعات است که در تحلیل \lr{SHAP} به‌صورت افتِ
اهمیت لگ‌های میانی خود را نشان می‌دهد.

الگوی رژیمی در افق روزانه نیز تکرار می‌شود: \lr{price\_D-1} از \lr{3.1109} به
\lr{5.1335} می‌رسد، یعنی حدود ۶۵ درصد افزایش، و \lr{exog\_1\_D0} از
\lr{3.4393} به \lr{3.0130} افت می‌کند. پس این یافته مختص افق ساعتی نیست.

\subsection*{یک روز، به‌تفصیل}

\begin{figure}[htbp]
  \centering
  \includegraphics[width=0.85\textwidth]{15_shap_waterfall_case_study}
  \caption{نمودار آبشاری برای یک روز منفرد: تجزیه‌ی فاصله‌ی میان پیش‌بینیِ آن
  روز و میانگین پیش‌بینی‌ها به سهم تک‌تک ویژگی‌ها.}
  \label{fig:shap-waterfall}
\end{figure}

تمام آنچه تا اینجا آمد، میانگین بر هزاران نمونه بود. \cref{fig:shap-waterfall}
همان تجزیه را برای یک روز منفرد نشان می‌دهد، و کارکردی متفاوت دارد.

میانگین‌ها به پرسش‌های پژوهشی پاسخ می‌دهند، اما در کاربرد عملی، آنچه به کار
می‌آید تبیینِ همان پیش‌بینیِ خاصی است که پیشِ روست. نمودار آبشاری نشان می‌دهد که
پیش‌بینیِ یک روزِ مشخص چگونه از میانگین فاصله گرفته و کدام ویژگی‌ها آن را به کدام
سو برده‌اند. همین قابلیت است که ابزار عملیاتیِ \cref{sec:5-3} بر آن بنا شده است:
عددی که با تبیین همراه باشد، برای بهره‌بردار قابل ارزیابی است، و عددی که تنها
یک عدد باشد، نیست.

\subsection*{آنچه این بخش ادعا نمی‌کند}

سه چیز از این تحلیل نتیجه نمی‌شود، و ننوشتنشان دعوت به بدخوانی است.

\textbf{این‌ها «عوامل تعیین‌کننده‌ی قیمت برق آلمان» نیستند.} این‌ها ویژگی‌هایی
هستند که \emph{یک مدل درختی}، آموخته بر \emph{یک مجموعه‌داده} و در
\emph{یک پنجره‌ی زمانی}، بیش از بقیه به آن‌ها تکیه کرده است. \lr{SHAP} خروجی مدل
را تجزیه می‌کند، نه بازار را.

\textbf{این‌ها به چهار مدل دیگر تعمیم نمی‌یابند.} \lr{LEAR-LASSO}، \lr{LSTM} و
\lr{SARIMAX} تحلیل نشده‌اند. اینکه دقت آن‌ها از \lr{LightGBM} تفکیک‌پذیر نیست، به
هیچ روی به معنای آن نیست که از اطلاعات یکسانی استفاده می‌کنند؛ ممکن است سه مسیر
متفاوت به دقتی مشابه رسیده باشند، و این خودْ پرسشی است که این پژوهش باز
می‌گذارد.

\textbf{هم‌بستگیِ میان ویژگی‌ها، سهم‌ها را میان اعضای یک خانواده جابه‌جا
می‌کند.} لگ‌های قیمت به‌شدت با هم هم‌بسته‌اند، و وقتی چنین است، توزیع اهمیت
میان \lr{price\_D-1} و \lr{price\_D-2} تا حدی به ساختار خودِ درخت وابسته است.
به همین دلیل تمام تفسیرهای این بخش در سطح \emph{خانواده‌ی ویژگی} بیان شده‌اند و
نه در سطح ستون‌های منفرد، و تمام مقایسه‌ها میان دو حالتِ همان خانواده انجام شده
است \lr{—} جایی که ساختار هم‌بستگی ثابت می‌ماند و اثرش تا حد زیادی حذف می‌شود.
